\documentclass[11pt]{article}
\usepackage{geometry}
\geometry{verbose}
\setlength{\parskip}{6pt}
\setlength{\parindent}{0pt}
\usepackage{color}
\usepackage{amsmath, amssymb}
\usepackage[authoryear]{natbib}
\usepackage[unicode=true,
 bookmarks=false,
 breaklinks=false,pdfborder={0 0 1},backref=section,colorlinks=false]
 {hyperref}

\makeatletter
\newcommand{\R}{\mathbb{R}}
\newcommand{\E}{\mathbb{E}}
\newcommand{\Cov}{\mathrm{Cov}}
\newcommand{\Var}{\mathrm{Var}}
\providecommand{\tabularnewline}{\\}
\usepackage{graphicx}
\usepackage{float}
\usepackage{url}
\usepackage{array}
\usepackage{enumitem}

\title{\Huge Assignment 4}
\author{\Large Chris Conlon}
\date{\Large Fall 2026 --- due at the start of Session 8}

\makeatother

\begin{document}

\title{Assignment 4}
\maketitle
\begin{center}
\begin{tabular*}{0.9\textwidth}{@{\extracolsep{\fill}}@{\extracolsep{\fill}}l@{\extracolsep{\fill}}l@{\extracolsep{\fill}}l}
Econometrics I & $\qquad$ & Professor Chris Conlon\tabularnewline
NYU Stern &  & Email: \href{mailto:cconlon@stern.nyu.edu}{cconlon@stern.nyu.edu}\tabularnewline
\end{tabular*}
\par\end{center}

\textbf{Instructions.} This is the first of the short biweekly assignments: two problems,
covering the instrumental-variables designs from Session 7 (examiner or ``judge'' instruments,
and shift-share exposure designs). \texttt{R} with \texttt{fixest} or Python with
\texttt{pyfixest}; \texttt{feols(y \~{} 1 | d \~{} z)} estimates 2SLS. Turn in a PDF with your
results and reasoning, and your code. Derive first, then simulate to check yourself.

\subsection*{1 Loan officers as instruments}

A bank receives $4{,}000$ small-business loan applications, each assigned \emph{at random}
to one of $40$ loan officers. Officer $j$ approves an application if the firm's
creditworthiness score clears the officer's personal cutoff $c_{j}$:
\begin{align*}
s_{i} &= u_{i} + v_{i}, \qquad D_{i} = \mathbf{1}\{s_{i} > c_{j(i)}\},
\qquad c_{j} \sim \mathcal{N}(0.8,\ 0.5^{2}) \text{ (one draw per officer)},
\end{align*}
where $u_{i} \sim \mathcal{N}(0,1)$ is firm quality (seen by the officer, not by you) and
$v_{i} \sim \mathcal{N}(0,1)$ is paperwork noise. Cutoffs are drawn once per officer.
The outcome is the firm's log revenue two years later:
\begin{align*}
y_{i} = 1 + \tau_{i} D_{i} + u_{i} + \varepsilon_{i}, \qquad
\tau_{i} = 1 - 0.4\, s_{i}, \qquad \varepsilon_{i} \sim \mathcal{N}(0, 0.5^{2}).
\end{align*}
The loan helps every firm, but helps firms with better scores less (they have other options).
The \textbf{average treatment effect} is $\E[\tau_{i}] = 1$.

\begin{enumerate}[label=(\alph*)]
\item Before simulating. OLS of $y$ on a binary $D$ is a difference of two group means
(Lecture 1's indicator trick): $\E[y \mid D=1] - \E[y \mid D=0]$. Substitute the outcome
equation and show this equals
\[
\underbrace{\E[\tau_{i} \mid D_{i}=1]}_{\text{effect on the treated}}
\;+\;
\underbrace{\E[u_{i} \mid D_{i}=1] - \E[u_{i} \mid D_{i}=0]}_{\text{selection}},
\]
the decomposition from Session 5. Then use the approval rule to sign each piece: which firms
get approved, what does that say about their $u_{i}$, and what does it say about their
$\tau_{i}$ relative to the average of $1$? Which piece do you expect to dominate? Give the
economic story in one sentence.
\item Simulate the data. Regress $y$ on $D$. Then compute the two pieces from (a) directly
(you know $\tau_{i}$ and $u_{i}$) and confirm they add up. Compare each to $1$.
\item Build the instrument: for each application, the assigned officer's approval rate
computed \emph{leaving out application $i$}. Call it $z_{i}$ (``leniency''). Why leave $i$
out? (One sentence; think about what $D_{i}$ itself contributes to the officer's rate.)
\item Estimate the first stage (regress $D$ on $z$) and the 2SLS estimate of the effect of a
loan. Report both. Interpret the first-stage coefficient in words.
\item State the three conditions this design needs --- relevance, exclusion, and
\emph{monotonicity} --- in the language of loan officers. Which one does random assignment
buy you, which one is about how officers behave, and which one is about whether an officer's
strictness affects revenue \emph{other than} through the loan?
\item Your 2SLS estimate should be noticeably below the ATE of $1$. Explain who the
\emph{compliers} are in this design (which firms' approval actually depends on the officer they
drew), use the model to say what their $s_{i}$ looks like, and hence why the LATE is below the
ATE here. Confirm by computing the average of $\tau_{i}$ among simulated firms whose $s_{i}$
lies between the strictest and most lenient officer cutoffs.
\item Managerial summary, two sentences: the bank is considering making every officer
as strict as its strictest one. Which of your numbers speaks to that policy, and which
number would a consultant quoting ``the average effect of a loan'' be describing?
\end{enumerate}

\subsection*{2 A shift-share instrument for local demand}

A retailer wants to know how much local employment growth raises its stores' sales. It has
$300$ cities and $8$ industries. Cities specialize: each city has one \emph{dominant} industry,
chosen at random, with pre-period employment share $s_{ck} = 0.5$; the other seven industries
split the rest equally ($0.5/7$ each). Over the study period each industry gets a
\emph{national} growth shock $g_{k} \sim \mathcal{N}(0, 1)$, drawn once per industry. City
employment growth and store sales growth are:
\begin{align*}
x_{c} &= \underbrace{\textstyle\sum_{k} s_{ck}\, g_{k}}_{B_{c}} + 0.5\, \eta_{c} + \nu_{c},
\qquad\qquad
y_{c} = 0.5\, x_{c} + \eta_{c} + \varepsilon_{c},
\end{align*}
where $\eta_{c} \sim \mathcal{N}(0,1)$ is a local shock (a new highway, a plant closing) that
moves both employment and sales, and $\nu_{c}, \varepsilon_{c} \sim \mathcal{N}(0, 0.5^{2})$.
The retailer also has last period's sales growth, $y^{0}_{c} = \eta_{c} + \xi_{c}$ with
$\xi_{c} \sim \mathcal{N}(0,1)$: a noisy look at the local shock. The true effect is $0.5$.

\begin{enumerate}[label=(\alph*)]
\item Simulate. Regress $y$ on $x$ by OLS. Which direction is the bias, and why?
\item Construct the \emph{shift-share} (Bartik) instrument $B_{c} = \sum_{k} s_{ck} g_{k}$:
each city's exposure to the national shocks, weighted by what it was already doing. Estimate
2SLS with $B_{c}$ as the instrument for $x_{c}$. Compare to $0.5$.
\item Two different assumptions make $B_{c}$ a valid instrument: (i) the \emph{shares} are
as good as randomly assigned with respect to $\eta_{c}$ (Goldsmith-Pinkham, Sorkin and Swift
2020), or (ii) the \emph{shocks} $g_{k}$ are as good as randomly assigned, whatever the
shares look like (Borusyak, Hull and Jaravel 2022). Which one holds in this simulation, and
what in the setup tells you?
\item A colleague proposes a check: regress the pre-period outcome $y^{0}_{c}$ on $B_{c}$. What
should you find if the instrument is valid, and what do you find?
\item Now break it. Suppose industry $1$ is booming ($g_{1} = 2$ instead of a random draw), and
that cities with a good local shock are more likely to be industry-$1$ cities: make industry $1$
the dominant industry with probability $\Phi(\eta_{c})$ (the standard normal CDF), otherwise
draw the dominant industry at random from the other seven. Redo (b) and (d). What happened to
the 2SLS estimate, which of the two assumptions in (c) failed, and did the check in (d) catch it?
\item Inference, one paragraph: there are $300$ cities but only $8$ industry shocks.
Explain why the usual city-clustered standard errors overstate how much independent variation
the instrument has, and what the ``number of observations'' effectively is in this design.
(You do not need to implement a correction; Ad\~ao, Koles\'ar and Morales 2019 do.)
\end{enumerate}

\end{document}
